\documentclass[12pt]{article}
\usepackage{amssymb, amstext, amsmath, amsthm}
\usepackage{hyperref}

\newcommand{\C}{\mathbb{C}}
\newcommand{\R}{\mathbb{R}}
\newcommand{\Q}{\mathbb{Q}}
\newcommand{\Z}{\mathbb{Z}}
\newcommand{\N}{\mathbb{N}}
\newcommand{\F}{\mathbb{F}}

\hypersetup{
    colorlinks,
    citecolor=black,
    filecolor=black,
    linkcolor=black,
    urlcolor=black
}

\newtheoremstyle{problem}  % <name>
        {10pt}                                               % <space above>
        {10pt}                                               % <space below>
        {\normalfont}                               % <body font>
        {}                                                  % <indent amount}
        {\bfseries}                 % <theorem head font>
        {\normalfont\bfseries:}         % <punctuation after theorem head>
        {.5em}                                          % <space after theorem head>
        {}                                                  % <theorem head spec (can be left empty, meaning `normal')>
\theoremstyle{problem}
\newtheorem{problem}{Problem}
\newtheorem{sol}{Solution}
\theoremstyle{remark}
\newtheorem{rmk}{Remark}

\begin{document}
\title{linear algebra pset}
\author{Edwin Lin}
\date{}
\maketitle
\newpage
\tableofcontents
\newpage

\section{Vector Spaces}
\begin{rmk}
We denote $\F$ as $\R$ or $\C$ unless otherwise specified.
\end{rmk}
\subsection{Basic Problems}
\begin{problem}
Prove that $-(-v) =v $ for every $v \in V$. 
\end{problem}
\begin{problem}
    Suppose $a \in \F, v \in V$ and $av = 0$. Prove that $a = 0$ or $v = 0$.
\end{problem}
\begin{problem}
    Suppose $v,w \in V$. Explain why there exists a unique $x \in V$ such that $v +3x = w$.
\end{problem}
\begin{problem}
    The empty set is not a vector space. The empty set fails to satisfy only one of the requirements listed in the definition of a vector space. Which one? 
\end{problem}

\begin{problem}
    Let $\infty$ and $-\infty$ denote two distinct objects, neither of which is in $\R$. Define an addition and scalar multiplication on $\R \cup\{\infty,-\infty\}$ as you could guess from the notation. Specifically, the sum and product of two real numbers is as usual, and for $t \in \R$ define \\ 
\begin{equation*}
t\infty = \begin{cases}
             -\infty  & \text{if } t < 0 \\
             0  & \text{if } t =  0 \\
             \infty & \text{if } t > 0
       \end{cases} \quad
t(-\infty) = \begin{cases}
             \infty  & \text{if } t < 0 \\
             0  & \text{if } t = 0 \\
             -\infty &\text{if } t >0
       \end{cases}
\end{equation*}
and
\begin{align*}
    t + \infty = \infty + t &= \infty + \infty = \infty \\
    t+ (-\infty) = (-\infty) + t &= (-\infty) + (-\infty) = -\infty \\
    \infty + (-\infty) &= (-\infty) + \infty = 0
\end{align*}
\end{problem}
\subsection{Subspaces}
\begin{problem}
For each of the following subsets of $\F^3$, determine whether it is a subspace of $\F^3$. 
\begin{enumerate}
    \item $\{(x_1,x_2,x_3) \in \F^3 \mid x_1 + 2x_2  + 3x_3 = 0\}$
    \item $\{(x_1,x_2,x_3) \in \F^3 \mid x_1 + 2x_2  + 3x_3 = 4\}$
    \item $\{(x_1,x_2,x_3) \in \F^3 \mid x_1x_2 x_3 = 0\}$
    \item $\{(x_1,x_2,x_3) \in \F^3 \mid x_1 =5x_3\}$
\end{enumerate}
\end{problem}

\begin{problem}
    Show that the set of differentiable real-valued functions $f$ on the interval $(-4,4)$ such that $f'(-1) = 3f(2)$ is a subspace of $\R^{(-4,4)}$.
\end{problem}

\begin{problem}
    Suppose $b \in \R$. Show that the set of continuous real-valued functions $f$ on the interval $\left[0,1\right]$ such that $\int_{0}^{1} f = b$ is a subspace of $\R^{\left[0,1\right]}$ if and only if $b = 0$.
\end{problem}

\begin{problem}
    Is $\R^2$ a subspace of the complex vector space $\C^2$?
\end{problem}

\begin{problem} Show the following.
    \begin{enumerate}
        \item Is $\{(a,b,c) \in \R^3 \mid a^3 =b^3\}$ a subspace of $\R^3$
        \item Is $\{(a,b,c) \in \R^3 \mid a^3 =b^3\}$ a subspace of $\C^3$
    \end{enumerate}
\end{problem}

\begin{problem}
    Prove or give a counterexample: If $U$ is a nonempty subset of $\R^2$ such that $U$ is closed under addition and under taking additive inverses (meaning $-u \in U$  whenever $u \in U$), then $U$ is not a subspace of $\R^2$.
\end{problem}

\begin{problem}
    Given an example of a nonempty subset $U$ of $\R^2$ such that $U$ is closed under scalar multiplication, but $U$ is not a subspace of $\R^2$.
\end{problem}


\newpage
\section{Solutions}
\subsection{Vector Spaces}
\subsubsection{Basic Problems}
\begin{sol}
    We show that $-(-v)$ is the additive inverse of $-v$ for every $v \in V$.
    \begin{align*}
        -v + (-(-v)) &= -v + (-1)(-v)\\
        &= (1 + (-1))(-v)  \\
        &= 0(-v)  \\
        &= 0
    \end{align*}
    Therefore, $-(-v)$ is the additive inverse of $-v$, hence $-(-v) = v$. 
\end{sol}
\begin{sol}
    Suppose $a \in \F, v \in V$, and $av =0$. Assume that $a \neq 0$ and $v \neq 0$. Then
    \begin{align*}
        av &= 0 \\
        a^{-1}av &= a^{-1}0 \\
        v &=  0
    \end{align*}
    But, $v \neq 0$. Thus, by contradiction, it must be that $a = 0$ or $v = 0$.
\end{sol}
\begin{sol}
    Suppose $v,w \in V$. First, we show existence. Consider
    \begin{align*}
        x = \frac{1}{3}\left(w-v\right) \in V
    \end{align*}
    We obtain $v +3x = w$. Next, we show uniqueness. Assume there exists $x,x' \in V$ such that $v + 3x = w =v +3x'$. Then
    \begin{align*}
        v + 3x &= v + 3x' \\
        3x &= 3x' \\
        x &= x'
    \end{align*}
    Therefore, there exists a unique $x \in V$ such that $v  + 3x = w$.
\end{sol}
\begin{sol}
    The additive identity property.
\end{sol}
\begin{sol}
    No, $\R \cup \{\infty, -\infty\}$ is not a vector space over $\R$. Consider $t \in \R\setminus\{0\}$. Then 
    \begin{align*}
        (t + \infty) + (-\infty) = \infty + (-\infty) = 0
    \end{align*}
    but
    \begin{align*}
        t + (\infty + (-\infty)) = t + 0 = t \neq 0
    \end{align*}
    Thus, it does not satisfies associativity.
\end{sol}
\subsubsection{Subspaces}
\begin{sol}
Let us denote each subsets of $\F^3$ as $U$. \\
\begin{enumerate}
    \item 
    \textit{\underline{Additive Identity?}} We have that $0+2(0) + 3(0) = 0$. Thus $0 \in U$. \\ \\
    \textit{\underline{Closed under addition?}} Consider $v,w \in U$ such that $v:=(v_1,v_2,v_3)$ and $w := (w_1, w_2,w_3)$. Since $v \in U$, then 
    \begin{align*}
        v_1 + 2v_2 + 3v_3 = 0
    \end{align*}
    Similarly, $w \in U$
    \begin{align*}
        w_1 + 2w_2 + 3w_3 = 0
    \end{align*}
    Adding them together, we get
    \begin{align*}
        (v_1 + w_1) + 2(v_2 + w_2) + 3(v_3 + w_3) = 0
    \end{align*}
    Therefore, $v + w \in U$. \\ \\
    \textit{\underline{Closed under scalar multiplication?}} Consider $\lambda \in \F$, $v \in U$. Then $\lambda v := (\lambda v_1, \lambda v_2, \lambda v_3)$ and 
    \begin{align*}
        \lambda v_1 + 2(\lambda v_2) + 3(\lambda v_3) &= \lambda(v_1 + 2v_2 + 3v_3)  \\ &= \lambda 0 \\ &= 0
    \end{align*}
    Therefore, $\lambda v \in U$. \\ \\
    Since $U$ satisfies the three conditions above, then $U$ is a subspace.
    \item 
    Not a subspace, we have that $0 + 2(0) + 3(0) = 0 \neq 4$. Hence $0 \not\in U$.
    \item
    Not a subspace, consider $(0,0,1), (1,1,0) \in U$. Then $(0,0,1) + (1,1,0) = (1,1,1)$ but the product of their components is non-zero. Thus $(1,1,1) \not\in U$.
    \item
    \textit{\underline{Additive Identity?}} Note that $0 = 5(0)$, thus $0 \in U$.  \\ \\
    \textit{\underline{Closed under addition?}} Let $v, w \in U$ and denote $v := (v_1,v_2,v_3)$ and $w := (w_1, w_2, w_3)$. Then 
    \begin{align*}
        v_1 + w_1 &= 5v_3 + 5w_3 \\
       &= 5( v_3 + w_3)
    \end{align*}
    Thus, $v + w  \in U$. \\ \\
    \textit{\underline{Closed under scalar multiplication?}} 
    Let $\lambda \in \F$ and $v \in U$, then for $\lambda v$ we have that 
    \begin{align*}
        \lambda v_1  &= \lambda(5v_3) \\
        &= 5(\lambda v_3)
    \end{align*}
    Thus, $\lambda v \in U$.
\end{enumerate}
\end{sol}
\begin{sol}
    Let us denote $D \subset \R^{(-4,4)}$ as the set of differentiable real-valued functions $f$ on the interval $(-4,4)$ such that $f'(-1) = 3f(2)$.\\ \\
    \textit{\underline{Additive Identity}}. Let us denote the additive identity function $0: (-4,4) \to \R$ such that $0(x) = 0$ for all $x \in (-4,4)$. Then we have that $0'(-1) = (3)0(2) = 0$. Thus $0 \in D$ \\ \\
    \textit{\underline{Closed under addition.}} Let $f,g \in D$, then 
    \begin{align*}
        (f+g)'(-1) &= f'(-1) + g'(-1) \\
        &= 3f(2) + 3g(2) \\
        &= 3(f(2) + g(2)) \\ 
        &= 3((f+g)(2))
    \end{align*}
    Thus, $f +g \in D$ \\ \\
    \textit{\underline{Closed under scalar multiplication.}} Let $\lambda \in \R$, $f \in D$. Then 
    \begin{align*}
        (\lambda f)'(-1) &= \lambda f'(-1) \\ &= 3(\lambda f)(2)
    \end{align*}
    Hence, $\lambda f \in D$. \\ \\
    Therefore, $D$ is a subspace of $\R^{(-4,4)}$.
\end{sol}
\begin{sol}
Suppose $b \in \R$. Let us denote $C \subset \R^{\left[0,1\right]}$ as the set of continuous real-valued functions $f$ on the interval $\left[0,1\right]$ such that $\int_{0}^{1} f  = b$. 
\\ \\
Suppose that $C$ is a subspace of $\R^{\left[0,1\right]}$. Let $f,g \in C$, then $f + g \in C$. So 
\begin{align*}
    \int_{0}^{1} f + g &= b \\
    \int_{0}^{1} f  + \int_{0}^{1} g  &= b \\
    2b &= b
\end{align*}
Hence, it must be that $b = 0$. Conversely, assume that $b = 0$.  \\ \\
\textit{\underline{Additive Identity?}} \\
Define the additive identity function $0: \left[0,1\right] \to \R$ such that $0(x) = 0$ for all $x \in \left[0,1\right]$. Then
\begin{align*}
    \int_{0}^10 = 0
\end{align*}
Thus, $0 \in C$.  \\ \\
\textit{\underline{Closed under addition?}} \\
Let $f,g \in C$ , then consider $f + g$
\begin{align*}
   \int_{0} ^{1} f + g &= \int_{0}^{1} f + \int_{0}^{1}g \\
   &= 0
\end{align*}
Hence, $f + g \in C$ \\ \\
\textit{\underline{Closed under scalar multiplication?}}
Let $\lambda \in \R$ and $f \in C$, then consider $\lambda f$
\begin{align*}
    \int_{0}^{1} \lambda f &= \lambda\int_{0}^{1}f \\
    &=\lambda(0)  \\
    &= 0
\end{align*}
We get that $\lambda f \in C$. \\ \\
Therefore $C$ is a subspace of $\R^{\left[0,1\right]}$. 
\end{sol}
\begin{sol}
    No, $\R^2$ is not a subspace of the complex vector space $\C^2$. Consider $i \in \C$ and $(0,1) \in \R^2$. We see that $i(0,1) \not \in \R^2$.
\end{sol}
\end{document}